\honest{Amazing Breakthrough Day: April 1st}{Eliezer Yudkowsky}{2008}

April 1st is already April Fools' Day, you may say. Yes, and that makes it ``the ideal cover'' for Amazing Breakthrough Day.

As I argued the day before, in ``The Beauty of Settled Science,'' science news covers only the frontier, where a new discovery is often controversial, rests on one experiment, is too hard to follow, and is later shown to be wrong. So people never see the solid, understandable science. \nb{Later evidence supports this; see ``The Beauty of Settled Science.''}

My proposal: on April 1st, journalists who care about science report old breakthroughs ``under the protective cover of April 1st.'' Sample headlines: ``BOATS EXPLAINED: Centuries-Old Problem Solved By Bathtub Nudist''; Königsberg tourists learn they cannot cross each bridge once; breathing turns out to be a kind of burning. Every one of them is true. They just did not happen yesterday. I do not say how readers, on the day they expect hoaxes, are to learn that these stories are true.

Old breakthroughs can be understood without a degree. Think of Archimedes's ``Eureka!'' when he understood why a ship floats; you can repeat the experiment in your bathtub. \nb{In the story, first written down by Vitruvius two centuries later, the bath and the ``Eureka!'' concern measuring the volume of King Hiero's crown, not ships. Archimedes did explain floating, in \textsc{On Floating Bodies}.} All of modern science is built on such discoveries, and reporting only the frontier is like walking into a film three-quarters of the way through.

If your editor says readers will not care, point out that Reddit and Digg readers vote up good explanations of old science. I give no example. Tell the editor that otherwise the paper will have to sell drugs to make payroll.

Then my larger idea: on the Internet, a good new explanation of old science ``is news and it spreads like news.'' Why shouldn't newspapers treat it that way? But ``all this is too visionary for a first step,'' so I settle for the holiday. \nb{The larger idea does not depend on April 1st.}

``April 1st. Put it on your calendar.''
